\subsection{Fidelity and the switched picture}
\label{sec:fidelity}

Near-baseline accuracy does not mean that $R$ steps of size $1/R$ reproduce the unit
step.
On ConvNeXt-T the refined endpoint ($R{=}100$) sits $\angleUntiedRefine^\circ$ from the
native one, with a relative distance of $\relDistUntied$ while the norms differ by
$3\%$; on the weight-tied model the gap is $\angleTiedRefine^\circ$ and $\relDistTied$
(\cref{app:solver-endpoints}; \cref{fig:refinement}, right).
The untied gap accumulates block after block, whereas the tied one stops growing after
two or three steps.
Stage~4 and the head nevertheless classify the refined state at a small cost.

The refined path is faithful to the native network in the statistics that matter for
transport.
On the untied stage, straightness and consecutive cosine at $R{=}1$ and $R{=}100$ agree
to within $\gramFidelityTol$ (\cref{tab:coherence}).
On the tied stage the order of the three regimes is preserved while the path
straightens ($\strTiedNative\!\to\!\strTied$), as a single autonomous field should under
refinement.
Together with the usable limit, this is the sense in which the probe is faithful: the
horizon is kept, the decoder still reads the state, and the Gram geometry of the
trained stage is recovered.

The same measurements give the switched picture.
The residual norm of an untied stage jumps at every block boundary
(\cref{fig:refinement}, right), and replacing the piecewise-constant schedule by a
bilinear blend of consecutive block weights costs $\bilinearDropLow$ to
$\bilinearDropHigh$ points on every architecture (\cref{app:bilinear}).
Each trained block is meant for its own unit of depth: a stage is a composition of one
ODE per block, not one ODE with weights smooth in depth.
